\documentclass[10pt,a4paper]{amsart}
\usepackage[margin=19mm]{geometry}
\usepackage{amsmath,amssymb,mathtools}
\usepackage[scr=rsfs,cal=euler]{mathalfa}
\usepackage{xfrac}
\usepackage[unicode,hidelinks,bookmarks=false]{hyperref}
\newcommand{\ie}{i.e.\ }
\newcommand{\cf}{cf.\ }
\newcommand{\resp}{respectively\ }
\newcommand{\Subsection}{\subsection}
\providecommand{\nobreakdash}{\nobreak\hbox{-}}
\setlength{\emergencystretch}{2em}
\multlinegap0pt
\usepackage{amssymb}
\newtheorem{theorem}{Theorem}
\def\mpn{\operatorname{mpn}}
\title{On the Existence of Minimal Parametric Networks}
\author{Arsen Kh. Galstyan}
\date{Source: arXiv:2607.23351v2 (CC BY 4.0)}
\begin{document}
\maketitle

\noindent\emph{Source and licence.} On the Existence of Minimal Parametric Networks. Version arXiv:2607.23351v2 (CC BY 4.0), distributed by arXiv under the Creative Commons Attribution 4.0 International licence, http://creativecommons.org/licenses/by/4.0/, as recorded in the arXiv licence metadata for this exact version and shown on the abstract page https://arxiv.org/abs/2607.23351v2. Full licence notice: \texttt{licenses/arxiv-2607.23351-CC-BY-4.0.txt}.\par\smallskip

\noindent\textit{Editorial scope.} Counted unit: the article's theorem thm:Ambrosio (source0.tex lines 634--636). The article's complete original proof of it is delivered with it (source0.tex lines 638--671, one \texttt{proof} environment of 34 lines). No mathematics is written, repaired or re-derived: the statement and the proof are the article's own lines, in the article's own notation, and no retained line is substituted or re-flowed.  No further source-local context is needed: every cross-reference of every retained line resolves inside the two delivered spans.  Nothing was omitted from any retained span, so the package carries no omission record.  No retained line contains a citation, so the wrapper carries no bibliography and no bibliography entry.  No retained line contains a figure environment, an image-inclusion command or drawing code, so no image is delivered and the package declares no asset dependency.  Editorial formatting only, in addition: an AMS \texttt{amsart} preamble that restates the article's own declarations and macros used by the retained lines, namely the environments \texttt{theorem} on separate counters, together with the standard packages those lines need.  The retained environments print in this document as Theorem 1 in order of appearance, whereas the article prints the same items with the numbering of its own full document.  The complete native source of the article as distributed in the arXiv e-print for this exact version is bundled unchanged as \texttt{sources/206020/source0.tex}.
\begin{theorem}\label{thm:Ambrosio}
Let $X$ be a compactly equipped space. Then, for any finite set $\mathcal{A}\subset X$ and for any connected graph $G = (V, E)$ with boundary $\mathcal{A}$, there exists a minimal parametric network of type $G$ joining $\mathcal{A}$.
\end{theorem}

\begin{proof}

Let $m = \mpn[G, \mathcal{A}]$. For each $n\in \mathbb{N}$ choose $g_n\in [G, \mathcal{A}]$ such that $|g_n| < m + \frac{1}{n}$. Then the sequence of lengths $|g_n|$ is bounded by $m+1$. 

Fix a point $a\in \mathcal{A}$. Since the graph $G = (V, E)$ is connected, for any vertex $v\in V$ there exists a path $\{v_1 = v, v_2\}, \{v_2, v_3\},\ldots , \{v_{k-1}, v_k = a\}$ connecting $v$ to $a$. By the triangle inequality we get
$$
\bigl|g_n(v)\, a\bigr| \le \sum_{i=1}^{k-1} \bigl|g_n(v_i)\, g_n(v_{i+1})\bigr| \le |g_n| \le m+1
$$
for all $n$. Hence, all points $g_n(v)$ lie in the closed ball $B = B_{m+1}(a)$. By the definition of a compactly equipped space, on $B$ there exists a compact topology $\tau$ such that the pseudometric of the space $X$ is lower semicontinuous on $B\times B$ with respect to the topology $\tau\times \tau$.

Let $V_{\text{int}} = V\setminus \mathcal{A}$ be the set of interior vertices. Any network $g\in [G, \mathcal{A}]$ with images in $B$ is uniquely determined by the restriction $g|_{V_{\text{int}}}$, since $g|_{\mathcal{A}} = \mathrm{id}$. Consider the space
$$
Y = \prod_{v\in V_{\text{int}}} B,
$$
endowed with the product topology $\tau_{Y}$ generated by $\tau$. By Tychonoff's theorem, $Y$ is compact.

Consider the sequence of image tuples $\bigl\{g_n(V_{\text{int}})\bigr\}\subset Y$. It is well known that in a compact space every net (in particular, every sequence) has a convergent subnet. Let $\bigl\{g_{\alpha}(V_{\text{int}})\bigr\}_{\alpha\in \Lambda}$ (provided by $g_{\alpha}(V_{\text{int}}) = g_{n(\alpha)}(V_{\text{int}})$) be a subnet converging in $\tau_Y$ to some point $\widetilde{g}(V_{\text{int}})\in Y$. This means that for each vertex $v\in V_{\text{int}}$, the net $\bigl\{g_{\alpha}(v)\bigr\}\subset B$ converges to $\widetilde{g}(v)\in B$ in the topology $\tau$. For boundary vertices $a\in \mathcal{A}$, set $\widetilde{g}(a) = a$. Thus, $\widetilde{g}\in [G, \mathcal{A}]$.

Estimate the length of $\widetilde{g}$. For any edge $\{u, v\}\in E$, we have $g_{\alpha}(u)\rightarrow \widetilde{g}(u)$ and $g_{\alpha}(v)\rightarrow \widetilde{g}(v)$ in the topology $\tau$. By the lower semicontinuity of the pseudometric of $X$ in the topology $\tau\times \tau$, we have
$$
\bigl|\widetilde{g}(u)\, \widetilde{g}(v)\bigr| \le \liminf_{\alpha} \bigl|g_{\alpha}(u)\, g_{\alpha}(v)\bigr|.
$$
Summing over all edges, we obtain
\begin{multline*}
|\widetilde{g}| = \sum_{\{u, v\}\in E} \bigl|\widetilde{g}(u)\, \widetilde{g}(v)\bigr| \le \sum_{\{u, v\}\in E} \liminf_{\alpha} \bigl|g_{\alpha}(u)\, g_{\alpha}(v)\bigr| \le \\ \le \liminf_{\alpha} \sum_{\{u, v\}\in E} \bigl|g_{\alpha}(u)\, g_{\alpha}(v)\bigr| = \liminf_{\alpha} |g_{\alpha}|.
\end{multline*}
Since $\bigl\{g_{\alpha}\bigr\}_{\alpha\in \Lambda}$ is a subnet of the sequence $\{g_n\}$, for which $\lim\limits_{n\rightarrow \infty} |g_n| = m$, we have $\lim\limits_{\alpha} |g_{\alpha}| = m$. Consequently, $\liminf\limits_{\alpha} |g_{\alpha}| = m$, whence
$$
\mpn[G, \mathcal{A}]\le |\widetilde{g}| \le m = \mpn[G, \mathcal{A}].
$$

Thus, $|\widetilde{g}| = \mpn[G, \mathcal{A}]$, which by definition means that $\widetilde{g}$ is a minimal parametric network of type $G$ joining $\mathcal{A}$.\qed

\end{proof}

\end{document}
